\documentclass[twocolumn]{aastex631}

\usepackage{amsmath}
\usepackage{multirow}
\usepackage{natbib}
\usepackage{graphicx} 
\usepackage{aas_macros}

\begin{document}

In this paper, we addressed the foundational question of whether the conjectured local symmetry that protects Degenerate Higher-Order Scalar-Tensor (DHOST) theories from classical ghost instabilities can survive in a more complete description that includes quantum corrections. The central problem is the potential conflict between the highly specific structure required by this symmetry and the generic form of leading-order quantum corrections, which are expected to arise in any effective field theory of gravity. To resolve this, we investigated the invariance of an effective action composed of a classically ghost-free DHOST Lagrangian supplemented by constant-coefficient Gauss-Bonnet and Weyl-squared terms.

Our methodology consisted of a direct, first-principles calculation of the variation of the complete effective action under the proposed local, field-dependent symmetry transformation. This transformation non-trivially mixes the scalar and metric sectors and is defined by two functions, $\Lambda(\phi, X)$ and $\alpha(\phi, X)$. Crucially, we made no prior assumption of invariance, instead computing the transformation of each term in the Lagrangian from the ground up to determine the precise conditions under which the action's variation would vanish.

The results of our analysis reveal a stark dichotomy. We have rigorously confirmed that the classical DHOST sector of the action is indeed invariant under the transformation. This invariance holds provided the defining functions $\Lambda$ and $\alpha$ satisfy a specific linear partial differential equation, $\Lambda + \alpha + 2X \partial_X \alpha = 0$. This finding provides direct evidence for the conjecture that the classical health of DHOST theories is a direct consequence of this underlying symmetry. In stark contrast, we demonstrated unequivocally that the constant-coefficient Gauss-Bonnet and Weyl-squared terms explicitly break this symmetry. The variation of these curvature-squared terms is non-zero and structurally independent of the variation of the classical terms, precluding any possibility of cancellation.

From these results, we have learned that there is a fundamental tension between the symmetry that guarantees classical viability and the standard form of leading-order quantum corrections. The unequivocal conclusion of this work is that a self-consistent, quantum-corrected DHOST theory cannot be constructed by naively appending generic, constant-coefficient curvature invariants to the classical Lagrangian. Such an approach destroys the very symmetry responsible for eliminating the pathological ghost degree of freedom, likely rendering the resulting theory physically unstable. This finding serves as a powerful new guiding principle for theoretical model building. It strongly suggests that for a DHOST theory to remain viable beyond the classical approximation, any quantum corrections must be incorporated in a more subtle, symmetry-preserving fashion. This could involve, for instance, promoting the coefficients of curvature-squared terms to be specific functions of the scalar field and its kinetic energy, or constructing entirely new classes of operators designed from the outset to be invariant. Ultimately, this work transforms a theoretical conjecture into a concrete and restrictive criterion, charting a more constrained and theoretically sound path toward a consistent theory of quantum-corrected modified gravity.

\end{document}
                